\subsection{Immediate-Switch Paired Intervals}
Table~\ref{tab:no-replay-ci} reports all requested paired bootstrap comparisons
for the five-seed immediate-switch experiment. Accuracy differences are percentage
points; ``before CO'' is a difference in proportions. Each interval uses
100,000 paired resamples of seeds 17, 23, 42, 101, and 202.

\begin{table}[h]
\centering
\caption{Immediate-switch minus the matched response: mean difference [95\% CI]}
\label{tab:no-replay-ci}
\scriptsize
\setlength{\tabcolsep}{3pt}
\begin{tabular}{lrr}
\toprule
Metric & Immediate-replay & Next epoch \\
\midrule
Before-CO proportion & $0.00[0.00,0.00]$ & $+0.40[0.00,+0.80]$ \\
Warn$\to$start (updates) & $0.0[0.0,0.0]$ & $-157.6[-252.0,-56.8]$ \\
Clean (points) & $-0.07[-0.16,0.00]$ & $-0.02[-0.14,+0.08]$ \\
PGD-10 (points) & $+0.22[-0.22,+0.66]$ & $+0.33[+0.08,+0.66]$ \\
Core time (min) & $+0.06[-0.18,+0.30]$ & $+0.59[+0.47,+0.71]$ \\
\bottomrule
\end{tabular}
\end{table}

\subsection{Offline Sensitivity Analysis}
Below we perform offline sensitivity analysis on saved held-out traces, without
retraining or modifying main-paper frozen parameters. The main protocol remains
$\tau=0.1663298875$, a 60-update near horizon, and the original CO definition.
The threshold scan is descriptive and is not used to replace the frozen value.
It uses the seven held-out CIFAR-10 event trajectories and three separately
reported randomized-FGSM non-event trajectories.

\subsubsection{Threshold Sensitivity}
Table~\ref{tab:tau-sensitivity} recomputes warnings from saved cosine traces.
Event recall, 60-update near recall, and non-event alarms are unchanged from
0.10 through 0.24. The median first-warning lead changes from 60 to 40 updates
at thresholds of 0.20 and above. Event and false-alarm counts are stable around
the frozen threshold; no post-hoc retuning is applied.
Here ``median first-warning lead'' uses the first active warning anywhere
before each event. Table~\ref{tab:transfer} instead reports ``median near lead'':
the earliest active warning restricted to the final 60 updates. At the frozen
threshold these deliberately different summaries are 60 and 40 updates,
respectively.

\begin{table}[h]
\centering
\caption{Offline threshold sensitivity on saved traces. Lead is the first-warning
lead; Table~\ref{tab:transfer} reports near lead.}
\label{tab:tau-sensitivity}
\scriptsize
\setlength{\tabcolsep}{3pt}
\begin{tabular}{rrrrr}
\toprule
$\tau$ & Recall & Near & Non-event alarms & Median first lead \\
\midrule
0.10 & 7/7 & 7/7 & 0 & 60 \\
0.12 & 7/7 & 7/7 & 0 & 60 \\
0.14 & 7/7 & 7/7 & 0 & 60 \\
0.16 & 7/7 & 7/7 & 0 & 60 \\
0.1663298875 (frozen) & 7/7 & 7/7 & 0 & 60 \\
0.18 & 7/7 & 7/7 & 0 & 60 \\
0.20 & 7/7 & 7/7 & 0 & 40 \\
0.22 & 7/7 & 7/7 & 0 & 40 \\
0.24 & 7/7 & 7/7 & 0 & 40 \\
\bottomrule
\end{tabular}
\end{table}

\subsubsection{Warning-Horizon Sensitivity}
At the frozen threshold, Table~\ref{tab:horizon-sensitivity} varies only the
reporting horizon. Six of seven events have an active warning within 20
updates, and all seven are covered by 40 updates. The registered 60-update
horizon is retained in the main paper.

\begin{table}[h]
\centering
\caption{Offline warning-horizon sensitivity}
\label{tab:horizon-sensitivity}
\scriptsize
\setlength{\tabcolsep}{6pt}
\begin{tabular}{rrrr}
\toprule
Horizon (updates) & Near events & Near recall & Non-event alarms \\
\midrule
20  & 6/7 & 85.7\% & 0 \\
40  & 7/7 & 100\% & 0 \\
60 (frozen) & 7/7 & 100\% & 0 \\
80  & 7/7 & 100\% & 0 \\
100 & 7/7 & 100\% & 0 \\
\bottomrule
\end{tabular}
\end{table}

\subsubsection{CO-Definition Sensitivity}
The original event requires a drop of at least 0.20 from the previous PGD-10
peak and accuracy at most 0.10. Table~\ref{tab:event-sensitivity} changes these
labeling constants only after training. All seven trajectories remain events
and are detected before the event under every variant. Near recall changes
because stricter labels place the registered event later or earlier relative
to saved warning windows; the main definition is unchanged.

\begin{table}[h]
\centering
\caption{Offline CO event-definition sensitivity}
\label{tab:event-sensitivity}
\scriptsize
\setlength{\tabcolsep}{3pt}
\begin{tabular}{lrrrrr}
\toprule
Definition & Drop & Ceiling & Events & Detected & Near \\
\midrule
Relaxed & 0.15 & 0.15 & 7 & 7/7 & 7/7 \\
Original & 0.20 & 0.10 & 7 & 7/7 & 7/7 \\
Higher ceiling & 0.20 & 0.15 & 7 & 7/7 & 7/7 \\
Strict & 0.25 & 0.05 & 7 & 7/7 & 5/7 \\
\bottomrule
\end{tabular}
\end{table}

\subsubsection{Long-Lead Signal Audit}
Table~\ref{tab:persistence} checks the two long-lead intervention traces. The
score is not continuously above threshold: only four post-warning windows are
active in each case. The first crossing is treated as a safety trigger after
observed geometry degradation, not as a calibrated estimate of time to collapse.

\begin{table}[h]
\centering
\caption{Signal persistence after long-lead warnings}
\label{tab:persistence}
\scriptsize
\setlength{\tabcolsep}{4pt}
\begin{tabular}{rrrrr}
\toprule
Seed & Lead & Windows & Active & Longest active run \\
\midrule
23  & 1568 & 80 & 4 (5.0\%) & 3 \\
202 & 1680 & 86 & 4 (4.7\%) & 2 \\
\bottomrule
\end{tabular}
\end{table}
